\section{VIV Object Format}
\label{sec:toolchain-object-format}

A VIV object file is the relocatable output produced by \texttt{viv32-asm} and
consumed by \texttt{viv32-ld}. Object files use the extension \texttt{.vo}.

A VIV object file contains three logical regions:

\begin{description}
    \item[Bytecode] The raw object-local bytes emitted by the assembler.
    \item[Symbol table] The object-local symbol definitions produced by labels.
    \item[Relocation table] The symbolic patches that must be applied by the linker.
\end{description}

The bytecode region is opaque. It may contain instructions, data, string literals,
zero padding, or any other emitted bytes. Relocation metadata is not embedded in the
bytecode stream. Instead, relocation records refer to byte offsets within the bytecode
region.

\subsection{Object-local addresses}

All offsets stored in a VIV object file are object-local byte offsets. The assembler
does not assign final executable addresses. Final addresses are assigned by the linker
after all input objects have been placed.

For example, a symbol recorded at object-local offset \texttt{0x00000100} does not
necessarily appear at final executable address \texttt{0x00000100}. If the linker
places the object at final base address \texttt{0x00000400}, the symbol's final
address is \texttt{0x00000500}.

\subsection{Object alignment}

The bytecode region of each VIV object is padded to a four-byte boundary. Padding bytes
are emitted as zeroes unless the assembler is padding executable text, in which case it
may emit \texttt{nop} instructions where appropriate.

This rule ensures that when \texttt{viv32-ld} places objects consecutively, the first
byte of each object begins at a word-aligned address. Individual labels and directives
remain the assembler's responsibility: instruction labels must be word-aligned, halfword
data must be halfword-aligned where required, and byte data requires no additional
alignment.

\subsection{File structure}

The initial VIV object format is designed to be simple and inspectable. Metadata is
stored before the raw bytecode blob. The bytecode blob is preceded by an explicit
length, so it may contain arbitrary byte values.

The logical structure is:

\begin{verbatim}
magic/version
symbol table
relocation table
bytecode length
bytecode bytes
\end{verbatim}

The exact textual form is:

\begin{verbatim}
VIVO 1

SYMBOLS <count>
<offset> <name>
...

RELOCATIONS <count>
<patch_offset> <symbol> <addend> <base> <sign> <value_shift> <width> <field_shift>
...

BYTES <byte_count>
<raw bytecode bytes follow>
\end{verbatim}

The \texttt{VIVO} header identifies the file as a VIV object file. The version number
allows future incompatible revisions of the object format.

\subsection{Symbol table}

The symbol table records labels defined by the assembled source file. Each symbol row
has the following form:

\begin{verbatim}
<offset> <name>
\end{verbatim}

The fields are:

\begin{description}
    \item[\texttt{offset}] The object-local byte offset at which the symbol is defined.
    \item[\texttt{name}] The object-level symbol name.
\end{description}

Symbol offsets are written as hexadecimal byte offsets.

By default, assembler labels are local to the object file. The assembler may mangle
local names so that they cannot collide with local names from other object files.
Symbols declared with \texttt{.nomangle} are recorded using their source-level name and
are visible to other object files during linking.

\subsection{Relocation table}

The relocation table records symbolic references that cannot be fully resolved by the
assembler. Each relocation row describes how the linker should compute a value and
insert that value into the bytecode blob.

Each relocation row has the following form:

\begin{verbatim}
<patch_offset> <symbol> <addend> <base> <sign> <value_shift> <width> <field_shift>
\end{verbatim}

The fields are:

\begin{description}
    \item[\texttt{patch\_offset}] Object-local byte offset of the encoded word or data
    item to patch.

    \item[\texttt{symbol}] Symbol whose final address is required.

    \item[\texttt{addend}] Signed adjustment added to the resolved symbol address before
    encoding.

    \item[\texttt{base}] Relocation base. This is either \texttt{abs} or \texttt{rel}.

    \item[\texttt{sign}] Field interpretation. This is either \texttt{u} for unsigned
    or \texttt{s} for signed.

    \item[\texttt{value\_shift}] Number of bits to shift the computed value right before
    insertion.

    \item[\texttt{width}] Width of the encoded field in bits.

    \item[\texttt{field\_shift}] Bit position of the encoded field within the patched
    word.
\end{description}

The relocation row is intentionally instruction-agnostic. The assembler understands
the instruction encoding and emits the appropriate relocation parameters. The linker
does not need to know which instruction or directive produced the relocation.

\subsection{Relocation calculation}

For an absolute relocation, the linker computes:

\begin{verbatim}
value = symbol_address + addend
\end{verbatim}

For a relative relocation, the linker computes:

\begin{verbatim}
value = symbol_address + addend - patch_address
\end{verbatim}

where \texttt{patch\_address} is the final executable address corresponding to
\texttt{patch\_offset}.

After computing the relocation value, the linker applies \texttt{value\_shift}:

\begin{verbatim}
encoded_value = value >> value_shift
\end{verbatim}

The linker then checks that the shifted value fits in \texttt{width} bits using the
specified signedness, masks the value to the requested width, shifts it left by
\texttt{field\_shift}, and inserts it into the patched word.

Alignment and semantic validity are the responsibility of the assembler. For example,
the assembler is responsible for ensuring that instruction labels are suitably aligned
and for choosing the appropriate relocation parameters for branch and jump
instructions.

\subsection{Relative relocation addends}

VIV-32 PC-relative branch and jump instructions encode offsets relative to the
instruction following the current instruction. The linker does not need to know this
architectural rule. Instead, the assembler records it using the relocation addend.

For example, a VIV-32 jump instruction at \texttt{patch\_address} requires:

\begin{verbatim}
target - (patch_address + 4)
\end{verbatim}

The assembler emits a relative relocation with an addend of \texttt{-4}:

\begin{verbatim}
target - patch_address - 4
\end{verbatim}

Thus the linker only applies the generic relative relocation rule.

\subsection{Example}

A simplified VIV object file may look like:

\begin{verbatim}
VIVO 1

SYMBOLS 2
00000080 _start
00000100 local_message

RELOCATIONS 2
00000000 _start       -4 rel s 2 26 0
00000088 local_message 0 abs u 0 15 8

BYTES 00000120
<raw bytecode bytes follow>
\end{verbatim}

The first relocation patches a PC-relative jump field. The assembler has encoded the
VIV-32 next-instruction-relative rule using the \texttt{-4} addend and has requested a
right shift of two bits for the word-scaled encoded offset.

The second relocation patches an absolute unsigned field with the final address of
\texttt{local\_message}. The linker does not need to know which instruction or
directive contains the field; it only applies the relocation record.
